% GENERATED FILE. Edit canonical lessons, then run npm run generate:books.
% canonical-source-hash: fnv1a64:3676f2e4ecf6af9b
% canonical-lessons: MR-C480-tabyet, MR-C480-laksan, MR-C480-thakva, MR-C480-asakta, MR-C480-malmalne

\chapter{Feeling Unwell}
\label{ch:mr-feeling-unwell}
\hlchaptermodality{480}

\begin{chapteropening}
\textbf{By the end of this chapter:} \emph{I can say how my health is and name a symptom.}

Complete the last lesson of chapter 480: I can say how my health is and name a symptom.
\end{chapteropening}

\section[\texorpdfstring{\mr{तब्येत} (tabyet)}{तब्येत (tabyet)}]{\mr{तब्येत} (tabyet) — one's health}

\label{lesson:MR-C480-tabyet}



\begin{quote}
Before the new one: say the Marathi for a cell, then the Marathi for the innermost heart.
\end{quote}

\subsection*{You'll want to know: \mr{तब्येत}}
\textbf{\mr{तब्येत}} — \emph{tabyet} — \textquotedblleft{}one's health\textquotedblright{}.

\begin{quote}
\mr{आजीची} \mr{तब्येत} \mr{आता} \mr{बरी} \mr{आहे}.
\end{quote}

\begin{quote}
\emph{ājīcī tabyet ātā barī āhe.}
\end{quote}

\begin{quote}
Grandmother's health is better now.
\end{quote}

Feminine.

\mr{तब्येत} \mr{बरी} \mr{नाही}, I am unwell; \mr{तब्येतीची} \mr{काळजी}, care of one's health.

\begin{cousinweb}
\textbf{\mr{तब्येत}} (\emph{tabyet}) is a loan from Persian \emph{tabiyat}, \textquotedblleft{}nature, health\textquotedblright{}, usually traced to Arabic.
\end{cousinweb}

\begin{grammarlens}[title={What you've built}]
One more word for describing what you lived and felt.
\end{grammarlens}

\subsection*{Your turn}
\begin{itemize}
\raggedright
  \item \emph{Say it:} \emph{tabyet}
  \item \emph{Say it:} \emph{tabyet}, once more
  \item \emph{Say it:} \emph{ājīcī tabyet ātā barī āhe.} — Grandmother's health is better now.
  \item \emph{Recall:} say the Marathi for a cell, then the Marathi for the innermost heart, then say \emph{tabyet} again
\end{itemize}

\begin{tcolorbox}[breakable,colback=teal!4,colframe=teal!35!black,title={Before you move on}]
What does \mr{तब्येत} mean? (\textquotedblleft{}One's health\textquotedblright{}.) Say it once more.
\end{tcolorbox}
\section[\texorpdfstring{\mr{लक्षण} (lakṣaṇ)}{लक्षण (lakṣaṇ)}]{\mr{लक्षण} (lakṣaṇ) — a symptom; a sign}

\label{lesson:MR-C480-laksan}



\begin{quote}
Before the new one: say the Marathi for the innermost heart, then the Marathi for one's health.
\end{quote}

\subsection*{You'll want to know: \mr{लक्षण}}
\textbf{\mr{लक्षण}} — \emph{lakṣaṇ} — \textquotedblleft{}a symptom; a sign\textquotedblright{}.

\begin{quote}
\mr{ताप} \mr{आणि} \mr{सर्दी} \mr{ही} \mr{कशाची} \mr{लक्षणं} \mr{आहेत}?
\end{quote}

\begin{quote}
\emph{tāp āṇi sardī hī kaśācī lakṣaṇaṁ āhet?}
\end{quote}

\begin{quote}
What are fever and a cold symptoms of?
\end{quote}

Neuter.

Neuter; plural \mr{लक्षणं}. A good sign is \mr{चांगलं} \mr{लक्षण}.

\begin{cousinweb}
\textbf{\mr{लक्षण}} (\emph{lakṣaṇ}) is a tatsama, taken from Sanskrit \emph{lakṣaṇa}, \textquotedblleft{}line, mark\textquotedblright{} (CDIAL 10882).
\end{cousinweb}

\begin{grammarlens}[title={What you've built}]
One more word for describing what you lived and felt.
\end{grammarlens}

\subsection*{Your turn}
\begin{itemize}
\raggedright
  \item \emph{Say it:} \emph{lakṣaṇ}
  \item \emph{Say it:} \emph{lakṣaṇ}, once more
  \item \emph{Say it:} \emph{tāp āṇi sardī hī kaśācī lakṣaṇaṁ āhet?} — What are fever and a cold symptoms of?
  \item \emph{Recall:} say the Marathi for the innermost heart, then the Marathi for one's health, then say \emph{lakṣaṇ} again
\end{itemize}

\begin{tcolorbox}[breakable,colback=teal!4,colframe=teal!35!black,title={Before you move on}]
What does \mr{लक्षण} mean? (\textquotedblleft{}A symptom; a sign\textquotedblright{}.) Say it once more.
\end{tcolorbox}
\section[\texorpdfstring{\mr{थकवा} (thakvā)}{थकवा (thakvā)}]{\mr{थकवा} (thakvā) — fatigue}

\label{lesson:MR-C480-thakva}



\begin{quote}
Before the new one: say the Marathi for one's health, then the Marathi for a symptom.
\end{quote}

\subsection*{You'll want to know: \mr{थकवा}}
\textbf{\mr{थकवा}} — \emph{thakvā} — \textquotedblleft{}fatigue\textquotedblright{}.

\begin{quote}
\mr{प्रवासानंतर} \mr{त्याला} \mr{खूप} \mr{थकवा} \mr{आला}.
\end{quote}

\begin{quote}
\emph{pravāsānantar tyālā khūp thakvā ālā.}
\end{quote}

\begin{quote}
He felt very tired after the journey.
\end{quote}

Masculine.

From \mr{थकणे}, to get tired: \mr{थकवा} \mr{येणे}, with the dative.

\begin{cousinweb}
\textbf{\mr{थकवा}} (\emph{thakvā}) is the Marathi noun of \mr{थकणे}, \textquotedblleft{}to get tired\textquotedblright{}, which goes back to the reconstructed \emph{*sthakk}, \textquotedblleft{}stop, halt\textquotedblright{} (CDIAL 13737).
\end{cousinweb}

\begin{grammarlens}[title={What you've built}]
One more word for describing what you lived and felt.
\end{grammarlens}

\subsection*{Your turn}
\begin{itemize}
\raggedright
  \item \emph{Say it:} \emph{thakvā}
  \item \emph{Say it:} \emph{thakvā}, once more
  \item \emph{Say it:} \emph{pravāsānantar tyālā khūp thakvā ālā.} — He felt very tired after the journey.
  \item \emph{Recall:} say the Marathi for one's health, then the Marathi for a symptom, then say \emph{thakvā} again
\end{itemize}

\begin{tcolorbox}[breakable,colback=teal!4,colframe=teal!35!black,title={Before you move on}]
What does \mr{थकवा} mean? (\textquotedblleft{}Fatigue\textquotedblright{}.) Say it once more.
\end{tcolorbox}
\section[\texorpdfstring{\mr{अशक्त} (aśakta)}{अशक्त (aśakta)}]{\mr{अशक्त} (aśakta) — weak}

\label{lesson:MR-C480-asakta}



\begin{quote}
Before the new one: say the Marathi for a symptom, then the Marathi for fatigue.
\end{quote}

\subsection*{You'll want to know: \mr{अशक्त}}
\textbf{\mr{अशक्त}} — \emph{aśakta} — \textquotedblleft{}weak\textquotedblright{}.

\begin{quote}
\mr{तापामुळे} \mr{ती} \mr{अशक्त} \mr{झाली} \mr{आहे}.
\end{quote}

\begin{quote}
\emph{tāpāmuḷe tī aśakta jhālī āhe.}
\end{quote}

\begin{quote}
The fever has left her weak.
\end{quote}

\mr{अ}-, not, + \mr{शक्त}, able; \mr{अशक्तपणा}, weakness. Unchanging.

\begin{cousinweb}
\textbf{\mr{अशक्त}} (\emph{aśakta}) is a tatsama, taken from Sanskrit \emph{aśakta}, \textquotedblleft{}unable, incompetent\textquotedblright{} (CDIAL 905).
\end{cousinweb}

\begin{grammarlens}[title={What you've built}]
One more word for describing what you lived and felt.
\end{grammarlens}

\subsection*{Your turn}
\begin{itemize}
\raggedright
  \item \emph{Say it:} \emph{aśakta}
  \item \emph{Say it:} \emph{aśakta}, once more
  \item \emph{Say it:} \emph{tāpāmuḷe tī aśakta jhālī āhe.} — The fever has left her weak.
  \item \emph{Recall:} say the Marathi for a symptom, then the Marathi for fatigue, then say \emph{aśakta} again
\end{itemize}

\begin{tcolorbox}[breakable,colback=teal!4,colframe=teal!35!black,title={Before you move on}]
What does \mr{अशक्त} mean? (\textquotedblleft{}Weak\textquotedblright{}.) Say it once more.
\end{tcolorbox}
\section[\texorpdfstring{\mr{मळमळणे} (maḷmaḷṇe)}{मळमळणे (maḷmaḷṇe)}]{\mr{मळमळणे} (maḷmaḷṇe) — to feel sick}

\label{lesson:MR-C480-malmalne}



\begin{quote}
Before the new one: say the Marathi for fatigue, then the Marathi for weak.
\end{quote}

\subsection*{You'll want to know: \mr{मळमळणे}}
\textbf{\mr{मळमळणे}} — \emph{maḷmaḷṇe} — \textquotedblleft{}to feel sick\textquotedblright{}.

\begin{quote}
\mr{बसच्या} \mr{प्रवासात} \mr{मला} \mr{मळमळतं}.
\end{quote}

\begin{quote}
\emph{bascyā pravāsāt malā maḷmaḷtaṁ.}
\end{quote}

\begin{quote}
I feel sick on bus journeys.
\end{quote}

Past (to him …): \mr{त्याला} \mr{मळमळलं}, \emph{tyālā maḷmaḷlaṁ}.

Dative and impersonal: \mr{मला} \mr{मळमळतं}, I feel queasy. \mr{मळमळ}, nausea.

\begin{cousinweb}
\textbf{\mr{मळमळणे}} (\emph{maḷmaḷṇe}) is a Marathi verb built on a doubled syllable; its origin is not settled.
\end{cousinweb}

\begin{grammarlens}[title={What you've built}]
One more word for describing what you lived and felt.
\end{grammarlens}

\subsection*{Your turn}
\begin{itemize}
\raggedright
  \item \emph{Say it:} \emph{maḷmaḷṇe}
  \item \emph{Say it:} \emph{maḷmaḷṇe}, once more
  \item \emph{Say it:} \emph{bascyā pravāsāt malā maḷmaḷtaṁ.} — I feel sick on bus journeys.
  \item \emph{Recall:} say the Marathi for fatigue, then the Marathi for weak, then say \emph{maḷmaḷṇe} again
\end{itemize}

\begin{tcolorbox}[breakable,colback=teal!4,colframe=teal!35!black,title={Before you move on}]
What does \mr{मळमळणे} mean? (\textquotedblleft{}To feel sick\textquotedblright{}.) Say it once more.
\end{tcolorbox}
